% part1_script.tex : 발표 대본(슬라이드별). presentation_script.tex 에서 \input 한다.
% 본문 수치는 deck/deck_spec_paper_report.json(notes.script, numbers)과
% docs/EXPERIMENT_PLAN_FINAL_BATCH_2026-10-04.md 8절, 두 추가 등록 문서, paper/claims/*/README.md 에서 옮겼다.
% 분량 기준: 약 300자/분(공백 제외). 선택 상자는 분량에 넣지 않는다. 글자 수는 count_part1.py 로 센다.

\slideheading{표지}{희소 관측 영구동토 지역의 활동층 두께 예측을 위한 물리 기반 기계학습의 라벨 수별 워크플로}{0.3}

안녕하십니까, 백승원입니다. 현장 관측값, 즉 라벨이 적은 영구동토 지역에서 활동층 두께를 예측할 때, 라벨 수에 따라 물리식과 기계학습을 어떻게 결합하고 평가할지를 정리한 결과를 보고드리겠습니다. 배경, 목적, 방법, 결과, 지도, 결론 순서입니다.

\pointat{바탕 지도(이 방법으로 만든 알래스카 1 km 지도), 제목의 ‘라벨 수별 워크플로’}

\slideheading{1}{활동층 두께와 관측 공백}{1.0}

먼저 대상입니다. 왼쪽 a처럼 영구동토 위 지표층은 여름에 녹았다가 겨울에 다시 업니다. 그 계절 최대 융해 깊이가 활동층 두께, ALT이고 단위는 cm입니다. 기반시설과 탄소·수문 과정에 영향을 주므로 넓은 영역의 지도가 필요합니다. 그런데 가운데 b처럼 관측은 점 관측망의 수백 지점에 몰려 있습니다. 직접 라벨은 1~km 셀로 묶어 17,467개이고 학습 원천 셀의 78에서 94\%가 알래스카입니다. 오른쪽 c의 네 지도는 같은 알래스카인데 영역 평균이 56에서 105~cm로 제품마다 다릅니다. 즉 대부분의 영구동토는 라벨이 거의 없는 새 지역이고, 어느 지도를 믿을지도 정해져 있지 않습니다.

\pointat{a 의 ‘융해 전선’과 ALT 화살표, b 의 출처 표, c 의 평균값 네 개}

\begin{optbox}
b 아래 수치처럼 지역별 Stefan 계수의 기하 평균은 레나델타 1.31에서 티베트 고원 11.03까지 다릅니다. 같은 열량에도 녹는 깊이가 지역마다 다르다는 뜻이고, 계수 재보정이 필요한 이유입니다.
\end{optbox}

\slideheading{2}{선행 연구와 남은 질문}{0.6}

ALT 기계학습 선행 연구 10편을 평가 방식으로 비교한 표입니다. 위 다섯 편은 학습 지역 안 무작위 교차검증, 아래 네 편은 거리 블록이나 지점 제외의 공간 교차검증입니다. 공통점이 셋입니다. 대상 지역을 통째로 빼고 평가한 연구가 없고, 라벨 수는 0개 또는 전량의 한 조건이며, 같은 라벨로 재보정한 물리 기준선이 없습니다. 맨 아래 주황색 행이 이 연구입니다. 남은 질문은 새 지역에서 라벨 수에 따라 무엇이 통하는가입니다.

\pointat{‘평가 유형’ 열, 맨 아래 주황색 행}

\begin{optbox}
조사 범위는 2019년부터 2026년까지의 영문 ALT 기계학습 논문이며 본문과 요약을 확인한 수준입니다. 이 범위에서 찾지 못했다는 뜻이지 ‘최초’라는 주장은 아닙니다.
\end{optbox}

\slideheading{3}{기존 연구의 한계와 이 연구의 개선}{0.7}

세 행이 평가 지역, 라벨 수, 물리 기준선입니다. 평가 지역은 지역 홀드아웃으로 바꿉니다. 대상 지역과 경계 100~km 완충을 학습에서 빼므로 학습에 없는 지역의 오차를 추정합니다. 라벨 수는 0개부터 전량까지의 곡선으로 평가해 몇 개부터 기계학습을 더할지 판단합니다. 물리 기준선은 같은 라벨로 재보정한 Stefan 식입니다. 물리식도 같은 라벨을 써야 재보정 몫과 기계학습 몫이 분리되기 때문입니다. 왼쪽 위 그래프처럼 같은 직접 ML이 무작위 분할에서는 15~cm 남짓, 지역 홀드아웃에서는 29~cm 가까이로 커집니다.

\pointat{왼쪽 위 그래프의 두 선, 가운데 열의 주황색 세 항목}

\slideheading{4}{연구 목적과 문제 정의}{0.7}

왼쪽 지도에서 원천 지역은 학습에 쓰는 지역, 대상 지역은 예측하려는 새 지역이고, 대상 지역과 100~km 완충 안의 셀은 학습에서 뺍니다. 원천 자료와 Stefan 식이 있고 대상 라벨이 $n$개일 때 무엇을 학습하고 무엇으로 평가하는가, 이것이 문제입니다. 오른쪽은 비교할 네 방법입니다. 원천 계수 Stefan은 대상 라벨을 쓰지 않고, 재보정 Stefan은 계수 $E$만 고칩니다. 물리 잔차 결합은 재보정 Stefan을 기본 답으로 두고 잔차를 학습하며, 직접 ML은 물리 결합 없이 배웁니다. 모두 같은 라벨과 같은 채점 블록을 씁니다.

\pointat{삽도의 100 km 완충 점선, 오른쪽의 라벨 수 축}

\slideheading{5}{연구 질문}{0.5}

질문은 라벨 수 구간별로 넷입니다. 라벨이 없을 때 물리 정보가 전이 오차를 줄이는가. 라벨 몇 개로 무엇을 먼저 고치는가. 라벨이 늘면 방법을 어떻게 고르는가. 어떻게 평가하고 지도는 무엇을 담는가. 오른쪽 열은 답을 한 줄로 미리 적은 것이고 근거는 결과에서 하나씩 보이겠습니다.

\pointat{오른쪽 ‘요지’ 열}

\slideheading{6}{용어}{0.6}

용어를 먼저 정합니다. 라벨은 현장에서 잰 ALT 관측값으로 1~km 셀 단위입니다. 원천 지역은 학습 지역, 대상 지역은 새 지역입니다. 원천 계수는 원천 라벨로 정한 Stefan 계수, 재보정 계수는 대상 라벨로 다시 맞춘 계수입니다. 물리 잔차 결합은 물리식이 남긴 오차를 기계학습이 보정하는 이 연구의 제안이고, 직접 ML은 물리 정보 없이 배웁니다. 판정은 셀 가중과 블록 등가중의 두 95\% 신뢰구간으로 하고, 차이가 ±0.5~cm 안이면 같다고 보는 것이 동등 범위입니다.

\pointat{‘물리 잔차 결합’ 행, ‘두 가중 신뢰구간’과 ‘동등 범위’ 행}

\slideheading{7}{사용 자료}{0.8}

자료는 라벨, 공변량, 물리 입력의 셋입니다. 가운데 표에서 주 4지역은 레나델타 점 관측 3,037행, 캐나다 750행, 러시아 서부와 동부의 CALM 지점 평균 31행과 30행입니다. 참조 지역 알래스카는 13,606행이지만 1~km 위치 343곳, 0.5° 블록 74개에 몰려 있습니다. 새 지역 러시아 중부와 티베트 고원은 57행과 132행입니다. 라벨 수 $n$은 행 수이고 행의 뜻이 지역마다 다릅니다. 공변량은 지형 6종, 기후 8종, 토양 9종, 위성 ALT 2종의 25종이며, 융해 도일은 ERA5-Land에서 계산합니다. 오른쪽 c처럼 모두 1~km 셀에 정렬하지만 정보 해상도는 토양 약 5~km, 기후 0.1°에 묶입니다.

\pointat{표의 알래스카 행(13,606 / 343 / 74), c 의 ‘정보 해상도’ 띠}

\begin{optbox}
라벨 연도는 1990년부터 2024년까지이고 예측은 2015년부터 2020년 기후값에 대한 정적 예측입니다. 연별 변동은 지금 공변량으로 예측되지 않아 다년 평균 규약을 씁니다.
\end{optbox}

\slideheading{8}{전체 연구 흐름}{0.6}

여섯 단계입니다. 자료를 1~km 셀로 정리하고, Stefan 식으로 원천 계수와 재보정 계수의 두 기준선을 만듭니다. 셋째, 기계학습을 물리 정보와 결합하는 방식들을 같은 학습기 CatBoost로 비교합니다. 넷째, 모든 방법을 지역 홀드아웃과 라벨 수 격자에서 같은 채점 블록으로 평가하며 판정 규칙은 실행 전에 등록했습니다. 다섯째, 결과를 라벨 수별 워크플로로 묶고, 여섯째, 알래스카 1~km 지도와 기존 제품 비교로 마무리합니다.

\pointat{단계 ①–⑥ 띠, 아래 주황색 되돌림 화살표(라벨을 더 모으면 앞 단계로)}

\slideheading{9}{Stefan 물리 모델의 뜻}{0.9}

Stefan 식은 여름에 지면으로 들어온 융해 열량으로 녹는 깊이를 구하는 열수지 해입니다. 열량은 영상 기온의 합인 융해 도일 TDD로 대신하고, 열전도도, 토양 수분, 밀도, 융해 잠열은 계수 $E$ 하나에 묶입니다. 그래서 ALT는 $E$ 곱하기 루트 TDD입니다. 왼쪽 아래는 러시아 서부 라벨 셀 하나인데, TDD 1,537도·일에 실측 118~cm를 원천 계수는 62~cm로 예측합니다. $E$가 지역마다 다르기 때문입니다. 오른쪽은 러시아 서부 라벨 31개에서 10개를 뽑아 계수를 다시 맞춘 예시로, 원천 계수 1.59가 2.11로 바뀝니다. 이것이 다음 쪽의 재보정입니다.

\pointat{가운데 식의 $E=\sqrt{2k/(\rho w L)}$, 오른쪽의 회색 선(1.59)과 파란 선(2.11)}

\begin{optbox}
$E$의 단위는 cm·(°C·일)$^{-1/2}$이고 주 지역의 원천 계수는 1.59에서 1.62로 거의 같습니다. 원천의 대부분이 알래스카 셀이기 때문입니다.
\end{optbox}

\slideheading{10}{Stefan 물리 모델과 계수 재보정}{0.8}

재보정 계수 $E_n$은 대상 라벨 $n$개의 최소제곱 계수 $E_{\mathrm{ls}}$와 원천 계수 $E_0$의 가중 평균이고 수축 강도 $\kappa$는 10입니다. 왜 섞느냐면, 라벨 몇 개로 계수를 통째로 바꾸면 오차가 오히려 커질 수 있기 때문입니다. 원천 계수에 라벨 10개의 무게를 주면 라벨 10개에서 두 계수의 중간이 되고 라벨이 늘수록 대상 계수에 가까워집니다. $\kappa$는 앞선 실험에서 정했고 3과 30은 민감도입니다. 그림은 레나델타 실제 라벨에서 10개를 뽑은 예시이며, 재보정 직선과 라벨의 차이인 잔차를 다음 쪽의 기계학습이 배웁니다.

\pointat{아래 식 $E_n=(nE_{\mathrm{ls}}+\kappa E_0)/(n+\kappa)$, 회색 원천 계수 선과 파란 재보정 선, 주황색 ‘잔차’}

\begin{optbox}
이 예시의 최소제곱 계수는 1.352, 재보정 계수는 1.486, 원천 계수는 1.619입니다. 라벨 10개이므로 재보정 계수가 정확히 중간입니다. 판정과 무관한 예시입니다.
\end{optbox}

\slideheading{11}{비교 방법과 모델 구조}{1.0}

비교한 방법은 물리 정보가 들어가는 위치로 나눈 다섯 구조입니다. 위 두 줄이 원천 계수와 재보정 Stefan이고, 주황색이 제안하는 물리 잔차 결합입니다. 재보정 Stefan 예측 $a$를 앵커로 두고 학습기 $g$는 라벨과 앵커의 차이인 잔차만 배웁니다. 예측은 $a$ 더하기 $\lambda$ 곱하기 $g(X)$이고 잔차 가중 $\lambda$는 0.25, 0.5, 1.0 가운데 고르며 전이 판정 기준은 0.25입니다. 배운 것이 없으면 예측이 물리식으로 돌아가는 구조입니다. 물리 입력 ML은 물리 예측을 입력 열로 넣는데 학습기가 그 열을 무시할 수 있습니다. 물리 유사라벨 증강은 Stefan 값을 학습 라벨로 더하고, 직접 ML은 물리 정보를 쓰지 않습니다. 모두 같은 CatBoost로 비교했고, 라벨 0개에서는 학습기 10종으로 결론이 학습기에 묶이지 않는지 확인했습니다.

\pointat{주황색 물리 잔차 결합 줄($\hat{y}=a+\lambda g(X)$), 오른쪽의 $\lambda$ 설명}

\begin{optbox}
CatBoost 설정은 반복 200회, 학습률 0.05, 깊이 3, seed 2개입니다. 학습기 10종은 CatBoost 3설정, 랜덤 포레스트, TabPFN v2, TabICL v2, 신경망 4종입니다. 공변량의 CCI ALT 2종은 같은 재분석 자료로 구동한 모형 출력이라 관측으로 취급하지 않습니다.
\end{optbox}

\slideheading{12}{물리 유사라벨 증강 대조 실험 설계}{0.6}

증강의 이득이 물리 정보에서 오는지, 학습 행이 늘어난 효과인지를 가르는 대조 실험입니다. 유사라벨의 자리와 수, 학습기, 채점 블록을 같게 두고 값만 바꾼 위약 넷을 만들었습니다. 순서를 섞은 유사라벨, 풀 평균 상수, 원천 평균 상수, TDD에 선형인 융해 지수입니다. 유사라벨은 라벨 블록 셀에 원천 행당 10개를 넣습니다. 오른쪽은 물리 유사라벨에서 위약을 뺀 오차 변화이고 음수가 물리 유사라벨이 나은 쪽입니다. 해석은 결과에서 하겠습니다.

\pointat{‘유사라벨 조건’ 다섯 줄, 오른쪽 그래프의 방향}

\slideheading{13}{평가 설계}{1.0}

핵심이므로 자세히 말씀드립니다. a는 지역 홀드아웃입니다. 대상 지역 셀과 100~km 완충을 학습 원천에서 뺍니다. b는 블록 절반 분할입니다. 0.5° 블록을 무작위로 절반씩 나누어 라벨은 노란 블록에서만 뽑고 채점은 초록 블록에서만 하며, 분할은 다섯 번입니다. c는 라벨 수 격자로 0, 3, 10, 40, 160, 320, 1000개와 전량이고 추출도 다섯 번씩입니다. d는 판정 규칙입니다. 두 가중의 신뢰구간이 모두 0 아래면 오차 감소, 모두 0 위면 오차 증가, 네 끝값이 ±0.5~cm 안이면 동등, 그 밖은 미결정입니다. 미결정은 차이를 확인하지 못했다는 뜻입니다. 맨 아래 줄처럼 가설과 판정 규칙을 커밋한 뒤 고정 seed로 실행하고, 재현 관문을 통과한 뒤에 봉인한 표를 열어 그 시각을 기록했습니다.

\pointat{b 의 노란 블록과 초록 블록, d 의 네 판정 그림, 맨 아래 ‘등록 → 실행 → 재현 관문 → 열람’}

\begin{optbox}
러시아 서부와 동부는 라벨 블록 셀이 15개 안팎이라 40개 이상은 레나델타와 캐나다 2지역 평균으로 봅니다. ±0.5~cm는 주 4지역 원천 계수 Stefan RMSE 21.7에서 43.0~cm의 약 1에서 2\%이며 결과 전에 고정한 규약입니다.
\end{optbox}

\slideheading{14}{검증 설계에 따른 오차}{0.7}

왜 지역 홀드아웃인지를 수치로 보입니다. 가로축은 시험 셀에서 가장 가까운 학습 셀까지의 거리, 세로축은 RMSE입니다. 보라색 파선인 직접 ML의 3지역 평균은 무작위 셀 분할 14.8~cm에서 지역 홀드아웃 28.7~cm로 거의 두 배가 되지만, 파란 실선인 재보정 Stefan은 20.5에서 22.2~cm로 거의 변하지 않습니다. kNNDM은 시험 셀과 학습 셀의 거리 분포를 지도 예측 상황에 맞춘 분할입니다. 무작위 분할은 학습 셀이 시험 셀 옆에 있어 새 지역 성능을 과대평가합니다.

\pointat{가로축의 거리 눈금, 두 선이 교차하는 지점}

\begin{optbox}
무작위에서 kNNDM으로 갈 때 직접 ML은 알래스카, 레나델타, 캐나다에서 5.8, 8.2, 14.2~cm, 재보정 Stefan은 0.3, 0.5, 4.3~cm 늘었습니다. 확률 표본이라면 무작위 검증도 타당하다는 반론이 있으나 이 연구의 라벨은 관측망에 몰린 자료라 확률 표본이 아닙니다.
\end{optbox}

\slideheading{15}{라벨 수별 워크플로 알고리즘}{0.5}

결과를 새 지역에 적용하는 절차입니다. 라벨이 없으면 원천 계수 Stefan을 쓰고 학습 범위 밖 셀을 표시합니다. 3개 이상이면 계수를 재보정하고 낮은 가중의 잔차를 더하며, 10개에서는 원천 계수 Stefan의 평균 절대 오차로 편향을 진단합니다. 40개 이상이면 라벨 블록 5겹 교차검증으로 후보 7개 가운데 하나를 고릅니다. 라벨 배치 효과는 지역에 따라 달라 사전 지정 절차는 무작위 배치로 확정되었습니다.

\pointat{흐름도의 네 상자, ‘교차검증 후보 7개’ 줄, ‘라벨 배치’ 줄}

\slideheading{16}{라벨 수 구간별 결과 요지}{0.5}

라벨 0개에서는 30개 대상 가운데 2~cm 넘게 나빠진 대상이 물리 유사라벨 증강 2개, 직접 ML 18개였습니다. 사후 서술입니다. 3에서 10개에서는 이득의 대부분이 계수 재보정이었습니다. 40에서 160개에서는 교차검증 선정이 고정 레시피에 뒤지지 않았고, 수백 개 이상인 알래스카에서는 물리 잔차 결합이 0.54~cm 작았으며 이득은 기후 격자 단위 보정이었습니다. 분산 배치 효과는 지역에 따라 달랐습니다.

\pointat{표의 다섯 행}

\slideheading{17}{물리 유사라벨의 정보 출처}{0.5}

위약 대조의 결과입니다. 라벨 0개에서 순서를 섞은 유사라벨보다 오차가 1.63~cm 작았고, 두 신뢰구간 모두 0 아래이며 보정 뒤 $p$는 0.001입니다. 다만 맨 아래 선형 융해 지수와의 차이는 0.48~cm로 0.5~cm 미만이고 라벨 10개에서는 동등입니다. 이득은 셀마다 다른 물리 값에서 오지만 대부분은 융해 도일에 선형인 정보로도 얻어집니다. 사전 등록 판정입니다.

\pointat{왼쪽이 물리 유사라벨이 나은 쪽, 회색 띠 ±0.5 cm. 첫 행의 점과 두 막대, 맨 아래 행이 띠 안에 있는 것}

\slideheading{18}{대상별 전이 오차 변화}{0.6}

같은 라벨 0개 조건의 대상별 지도입니다. 원 하나가 대상 하나이고 파랑이 오차 감소입니다. 왼쪽 직접 ML은 알래스카 삽도를 비롯해 어두운 원이 많고, 오른쪽 증강은 대부분 흰색에 가깝습니다. 2~cm 넘게 나빠진 대상이 18개 대 2개입니다. 예외는 왼쪽 아래 티베트 고원입니다. 원천 계수 Stefan 오차가 약 245~cm인 심부 레짐이라 직접 ML이 60~cm 가까이 낫습니다. 계수가 크게 틀린 곳에서는 물리식이 출발점이 될 수 없으므로 티베트는 따로 보고합니다.

\pointat{알래스카 삽도의 어두운 원, 티베트 삽도의 짙은 파란 원, 컬러바 방향}

\begin{optbox}
두 신뢰구간이 모두 0을 넘은 대상으로 세면 14개 대 8개입니다. 알래스카 대상에서는 증강도 원천 계수 Stefan보다 1.79~cm 나쁘지만, 직접 ML의 35.3~cm를 16.3~cm로 낮췄습니다. 손해를 줄이는 것이지 없애는 것은 아닙니다.
\end{optbox}

\slideheading{19}{기존 ALT 지도와 물리 기준선}{0.5}

라벨 0개의 비교 기준을 기존 지도까지 넓혔습니다. 원천 계수 Stefan 대비 라벨 연도의 융해 도일을 쓴 연도 정합 Stefan은 1.00~cm 작았고, 기존 지도 CCI v5는 학습 지점 주변 5~km를 가린 뒤에도 23.8~cm 컸습니다. 맨 아래처럼 직접 ML은 연도 정합 Stefan보다 3.24~cm 컸습니다. 따라서 라벨이 없을 때의 출발점은 물리식입니다. 기존 지도 비교는 결과를 본 뒤 설계했고 원인은 단정하지 않습니다.

\pointat{끊어진 축의 CCI v5 행, 맨 아래 직접 ML 행}

\begin{optbox}
Stefan·위성 제품 앙상블은 두 가중의 판정이 달라 미결정입니다. Wei 2026 지도는 2지역에서 미결정이었고, 대상 지역 안 CALM 지점도 학습에 썼으므로 우열을 말하지 않습니다.
\end{optbox}

\slideheading{20}{라벨 10개 계수 재보정}{0.7}

세로축은 원천 계수 Stefan 대비 오차 변화입니다. 왼쪽 주 4지역 평균에서 파란 재보정 Stefan은 10개에서 2.45~cm 줄였습니다. 두 신뢰구간 모두 0 아래, 보정 뒤 $p$ 0.001의 사전 등록 판정입니다. 그런데 오른쪽 레나델타·캐나다 평균에서는 재보정 Stefan이 0선 위로 올라갑니다. 유의하게 준 지역은 러시아 서부 하나로 11.96~cm이고 캐나다는 2.26~cm 나빠졌습니다. 캐나다는 지역 평균 계수가 원천과 거의 같지만 라벨 위치마다 계수가 달라, 뽑힌 라벨로 맞춘 계수가 채점 블록에 맞지 않았다고 해석합니다. 사후 해석입니다.

\pointat{왼쪽 파란 선의 라벨 10개 점(−2.45), 오른쪽 파란 선이 0선 위로 올라가는 것}

\begin{optbox}
오른쪽의 직접 ML은 라벨 40개부터 0선 아래로 내려옵니다. 라벨이 수십 개를 넘으면 직접 ML도 쓸 자료가 생기므로 뒤에서 방법을 라벨로 고르는 규칙을 둡니다.
\end{optbox}

\slideheading{21}{물리 잔차 결합의 추가 이득}{0.5}

주황색이 재보정 Stefan 위에 잔차 가중 0.25의 잔차 학습을 얹은 물리 잔차 결합입니다. 라벨 10개의 재보정 대비 추가 이득은 0.18~cm였고, 보정 전 $p$ 0.034, 보정 뒤 0.136이라 차이를 확인하지 못했습니다. 전량에서는 0.40~cm 작았지만 0.5~cm 미만입니다. 따라서 라벨 10개까지 이득의 대부분은 재보정입니다. 사전 등록 판정입니다.

\pointat{왼쪽의 파란 선과 주황 선이 거의 겹치는 것, 오른쪽의 두 띠가 겹치는 것}

\begin{optbox}
표형 파운데이션 모델을 잔차 학습기로 쓰면 재보정 대비 라벨 3개에서 0.32, 10개에서 0.55~cm의 작은 추가 이득이 있었습니다. 잔차 구조의 이득은 학습기 종류에 묶이지 않습니다.
\end{optbox}

\slideheading{22}{대상 라벨 교차검증 선정}{0.6}

라벨이 수십에서 수백 개일 때입니다. 라벨 안 블록 5겹 교차검증으로 후보 7개 가운데 오차가 가장 작은 방법을 고르는 선정 규칙을 고정 잔차 레시피와 비교했습니다. 레나델타·캐나다 평균에서는 40개에서 0.87, 160개에서 1.32~cm 작았고 두 행 모두 0.5~cm 한계 안에서 비열등입니다. 다만 이득은 캐나다에서 왔습니다. 네모인 캐나다는 2.10과 2.72~cm 줄었고 마름모인 레나델타는 +0.35와 +0.07~cm입니다. 사후 분석입니다.

\pointat{아래 두 행의 점과 막대, 캐나다 네모와 레나델타 마름모}

\begin{optbox}
주 4지역의 10개와 전량은 미결정이고 10개에서는 비열등이 성립하지 않습니다. 40개의 0.87~cm는 보조 신뢰구간에서 판정이 달라 분할 독립 가정에 의존합니다. 후보 7개는 Stefan 3종, 잔차 가중 2종, 증강 결합, 물리 유사라벨 증강이며 10개 미만이면 재보정 Stefan을 씁니다.
\end{optbox}

\slideheading{23}{워크플로 순차 적용 결과}{0.6}

배치, 라벨 10개 진단, 방법 규칙을 순서대로 적용한 결과입니다. 얕은 레짐 5지역을 새로 나눈 분할에서 평가했습니다. 위 묶음은 무작위 배치에 고정 레시피를 쓴 경우 대비로 40개 미결정, 160개 0.93~cm 감소입니다. 아래 묶음은 재보정 Stefan 대비로 10개 동등, 40개와 160개에서 1.35와 1.51~cm 감소입니다. 160개 이득은 캐나다에서 왔고, 러시아 동부, 알래스카, 레나델타는 일부 라벨 수에서 0.5~cm 넘게 나빴습니다. 배치 단계는 무작위로 귀결됐고, 같은 지역을 다시 나눈 것이라 독립 검증이 아닙니다.

\pointat{위 묶음 두 행과 아래 묶음 세 행, 오른쪽의 옅은 지역 기호}

\slideheading{24}{지역 내 라벨 수 곡선}{0.5}

라벨이 많은 지역 안의 결과입니다. 블록 절반을 채점에 남기고 분할을 25회 반복했습니다. 세로축은 재보정 Stefan 대비입니다. 알래스카의 주황색 물리 잔차 결합은 라벨 1000개에서 0.54~cm 작았고 두 신뢰구간 모두 0 아래입니다. 보라색 직접 ML은 1000개에서도 0선 위입니다. 캐나다와 레나델타는 차이를 확인하지 못했습니다. 점선은 오차 하한으로 알래스카 11.31~cm입니다. 첫 시험이 기각된 뒤 설계를 바꾼 시험이라 사후 분석입니다.

\pointat{알래스카 패널의 라벨 1000개 주황 점, ‘Error floor’ 점선}

\begin{optbox}
RMSE로는 14.41에서 13.87~cm이고 오차 하한은 공변량이 같은 셀 사이의 실측 차이로 추정했습니다. RMSE 개선률 3.7\%는 하한 위의 줄일 수 있는 오차 제곱으로는 19\% 감소이고 대회 설정의 39\%와 같은 방향입니다. 부록 10에 있습니다.
\end{optbox}

\slideheading{25}{보정량의 공간 분포}{0.6}

알래스카에서 기계학습이 고친 오차를 공간 규모로 나눴습니다. 오른쪽은 오차를 ERA5 기후 격자 사이 성분과 격자 안 성분으로 나눈 것입니다. Stefan 예측은 격자 안에서 같은 값이므로 나눌 수 있습니다. 총 0.52~cm 감소 가운데 격자 사이는 1.03~cm 감소, 격자 안은 0.01~cm로 동등입니다. 오차 제곱 감소의 약 99\%가 기후 격자 단위 편향 보정입니다. Stefan 오차 제곱의 72\%가 격자 안 변동인데 기계학습이 설명한 몫은 1\% 미만입니다. 즉 격자 안 상세도는 더하지 못했습니다. 알래스카 지역 안에 한정한 사후 분석입니다.

\pointat{왼쪽 지도(블록별 오차 변화, 파랑이 감소), 오른쪽 ‘Between’ 열과 ‘Within’ 열의 알래스카 점}

\begin{optbox}
세 지역 모두에서 격자 안 설명 비율은 8\% 미만이고, 세 지역 평균에서 격자 안 예측은 오차를 0.09~cm 오히려 늘렸습니다. 격자 묶음을 토양 융해 도일로 바꾼 민감도에서도 알래스카 결과는 같았습니다.
\end{optbox}

\slideheading{26}{라벨 위치 전략}{0.5}

다음 라벨을 어디에 둘지입니다. 캐나다에서 라벨 100개를 셀 무작위로 뽑은 경우와 공변량 공간에서 덜 덮인 곳을 차례로 고른 경우입니다. 무작위는 서쪽 큰 블록에 몰리고 공변량 분산은 북쪽 섬까지 퍼집니다. 캐나다 지역 안 50에서 100개에서 분산 배치는 무작위보다 2.54에서 3.21~cm, 새 지역 적용의 40개에서 2.39~cm 작았습니다. 그러나 레나델타에서는 블록 층화가 100개에서 약 1.7~cm 키웠습니다. 효과가 지역에 따라 다르므로 사전 지정 절차는 무작위로 확정했습니다.

\pointat{왼쪽의 큰 회색 원 하나, 오른쪽의 북쪽 섬 주황 점들}

\begin{optbox}
레나델타는 라벨이 몇 블록에 몰려 있어 블록마다 고르게 뽑으면 계수 추정이 치우친다고 보지만 원인 분석은 하지 않았습니다. 예측 분산 기반 능동 선택은 알래스카와 레나델타에서 무작위보다 0.12에서 2.79~cm 나빴습니다.
\end{optbox}

\slideheading{27}{독립 지역}{0.5}

학습에 쓰지 않은 지역에서도 결론이 유지되는지입니다. 러시아 중부를 더한 얕은 레짐 5지역에서, 라벨 40개까지 직접 ML이 원천 계수 Stefan보다 유의하게 나은 지역은 0개였습니다. 오른쪽 위 두 행은 전량에서 물리 잔차 결합이 원천 계수 Stefan보다 5지역 평균 3.19~cm 작았다는 것이고, 아래 두 행은 라벨 10개의 재보정 대비 추가 이득으로 동등입니다. 티베트는 심부 레짐이라 따로 보고합니다. 확인적 독립 지역이 4에서 5개라는 것이 한계입니다.

\pointat{지도의 흰 원, 오른쪽 ‘러시아 중부’ 행의 긴 블록 등가중 막대(미결정)}

\slideheading{28}{예측 구간}{0.4}

보조 결과인 예측 구간입니다. 점선이 목표 0.90입니다. 알래스카에서 보정한 보통의 기계학습 구간을 새 지역에 그대로 쓰면 포함률이 0.32에서 0.73으로 떨어지고, 지역을 집단으로 보정한 계층 conformal 구간은 주 4지역 평균 0.86입니다. 보정 집단이 4개라 유한표본 보장은 없으므로, 새 지역에는 지역 단위 보정이 필요하다는 것까지가 결론입니다.

\pointat{‘알래스카 보정 구간’ 행의 긴 범위 선, ‘지역 단위 보정’ 행의 큰 점}

\slideheading{29}{알래스카 1 km ALT 지도}{0.7}

지도입니다. 라벨이 많은 지역의 최종 방법인 물리 잔차 결합으로 알래스카 영구동토 구역 표시 셀 858,517개를 그린 1~km 지도이고 250~m 음영 지형 위에 얹었습니다. 정확도는 0.5° 블록 5겹 교차검증 값만 말씀드립니다. 물리 잔차 결합 13.80~cm, 재보정 Stefan 14.30~cm이고 90\% 예측 구간은 상수 폭 ±20.2~cm입니다. 세부 무늬는 지형 입력에서 오지만 정보 해상도는 기후 입력의 약 10~km에 묶입니다. 지도 셀의 73\%는 공변량이 학습 범위 밖이라 외삽으로 읽어야 합니다.

\pointat{북쪽 해안의 옅은 띠와 내륙의 짙은 띠, 컬러바 40–70 cm}

\begin{optbox}
외삽 셀에서는 구간의 포함률이 보장되지 않습니다. 범위 밖 공변량은 주로 지형 거칠기, 경사, TPI이며 라벨 343곳의 지형 범위가 전체 영역을 덮지 못하기 때문입니다. 세부 패널은 부록 12에 있습니다.
\end{optbox}

\slideheading{30}{라벨 수에 따른 지도 변화}{0.7}

같은 지역의 1~km 지도가 라벨 수에 따라 어떻게 바뀌는지를 레나델타로 보입니다. a는 라벨 0개의 원천 계수 Stefan입니다. b는 라벨 10개의 재보정 Stefan인데 계수가 1.62에서 2.04로 커져 전체가 깊어집니다. c와 d는 40개와 160개의 물리 잔차 결합으로, 잔차 가중 1.0이 선택되어 북동부가 뽑힌 라벨에 따라 크게 달라집니다. e의 전량 3,037개에서 안정됩니다. 요지는 라벨이 적을 때 지도는 어느 라벨이 뽑혔는지에 좌우되고 수백 개부터 안정된다는 것입니다. 정확도는 말하지 않으며, 라벨 40개 지도가 정확하다는 뜻이 아닙니다.

\pointat{윗줄 ALT, 아랫줄 전량 대비 변화. b 의 짙어진 색, c·d 의 북동부 짙은 영역, j 의 거의 흰 지도}

\slideheading{31}{표시 해상도와 정보 해상도}{0.6}

1~km 지도가 1~km만큼의 정보를 담는지 확인했습니다. 같은 지도를 왼쪽은 1~km 셀 값, 오른쪽은 0.1° 셀 평균으로 그렸습니다. 0.1° 평균이 1~km 지도 분산의 알래스카 97\%, 레나델타 61\%를 설명했습니다. 그래서 1~km는 표시 해상도이고 정보 해상도는 기후 입력 ERA5-Land 0.1°, 약 10~km와 토양 입력 약 5~km에 묶입니다. 앞의 이득 분해와 같은 방향이며, 그래서 고해상 ALT 지도라고 말하지 않습니다.

\pointat{a 와 b 의 같은 띠 구조, c 와 d 의 차이}

\begin{optbox}
0.1° 셀 안 1~km 값의 표준편차는 알래스카 1.8~cm(전체 9.5), 레나델타 1.1~cm(전체 1.8)입니다. 알래스카의 ERA5-Land 0.1° 격자는 약 4–6 × 11~km입니다.
\end{optbox}

\slideheading{32}{기존 ALT 제품과의 차이}{0.5}

우리 지도와 기존 제품이 어디서 얼마나 다른지입니다. 정확도 비교가 아니며 정확도는 같은 채점 셀 비교로만 판단했습니다. 영역 평균은 물리 잔차 결합 56.1~cm이고 CCI v5 105.5, Wei 2026 70.8, Aalto 2018 78.8, Yi·Kimball 105.1~cm로 기존 제품이 모두 깊습니다. 공간 상관은 재보정 Stefan과 0.96, 기존 제품과는 −0.15에서 0.46입니다. 차이가 큰 곳은 알래스카 내륙 등이고 원인은 단정하지 않습니다.

\pointat{a 의 mean 56.1 과 b–e 의 mean·r 값, g 의 올리브색 내륙}

\slideheading{33}{레나델타 입체 보기}{0.5}

레나델타 1~km 지도를 250~m 지형 위에 얹어 남동쪽에서 본 모습입니다. 지형은 열두 배 과장했고 남서쪽 구릉의 최고 높이는 약 580~m, 델타 평원은 약 30~m 아래이며 색 범위는 35에서 50~cm입니다. 평원의 ALT 무늬는 지형의 수로나 호수를 따르지 않고 기후 격자 입력의 띠를 따릅니다. 앞 쪽의 정보 해상도 결론을 눈으로 확인하는 그림이고, 축이 없으므로 모양과 상대 기복만 읽습니다.

\pointat{왼쪽 아래 구릉의 기복, 평원을 가로지르는 띠 무늬}

\slideheading{34}{결과 종합}{0.9}

결과를 라벨 수 구간별로 종합합니다. 왼쪽 주 4지역 평균에서는 재보정과 저가중 잔차 결합이 함께 내려가 10개에서 약 2.5~cm 줄었고 직접 ML은 0선 위입니다. 오른쪽 레나델타·캐나다 평균에서는 캐나다의 재보정 오차 때문에 라벨을 쓰는 방법들이 0선 위에 놓이지만, 짙은 갈색 점인 교차검증 선정은 40개부터 그 아래로 내려옵니다. 라벨이 없으면 물리식이며 기존 지도 제품보다 23.8~cm 정확했습니다. 10개면 재보정만으로 2.45~cm를 줄였고 잔차 학습의 추가 이득은 확인되지 않았습니다. 40개 이상에서는 교차검증 선정과 잔차 결합이 재보정 Stefan보다 작았고 160개에서 1.5~cm였습니다. 이것이 새 지역에서 학습, 평가, 추가 관측의 순서를 정하는 워크플로의 근거입니다.

\pointat{a 의 라벨 10개 점, b 의 짙은 갈색 점 세 개, 아래 권고 띠}

\begin{optbox}
160개의 1.5~cm는 같은 지역을 다시 나눈 풀 평균입니다. b 풀에서 라벨을 쓰는 방법이 0선 위에 놓이는 것은 캐나다의 재보정이 원천 계수보다 나빴기 때문이고, 그래서 40개부터는 고정 레시피가 아니라 대상 라벨로 고르는 규칙을 권고합니다.
\end{optbox}

\slideheading{35}{연구 의의와 차별성}{0.6}

의의는 셋입니다. 첫째, 지역 홀드아웃과 100~km 완충, 라벨 0개부터 전량까지의 곡선, 같은 라벨로 재보정한 물리 기준선을 함께 갖춘 ALT 연구는 조사 범위에서 찾지 못했습니다. 방법 하나하나는 새롭지 않고 기여는 평가 설계와 음성 결과를 포함한 결과, 워크플로입니다. 둘째, 무작위 분할은 직접 ML의 새 지역 오차를 절반으로 과소평가합니다. 셋째, 규모입니다. 평가 지역 7곳, 전이 대상 27개, 라벨 수 8단계, 방법 12종과 학습기 10종, 실행 전 등록한 가설 174개, 모형 적합 최소 863,229건입니다.

\pointat{왼쪽 표의 ‘등록 가설 174개’와 ‘모형 적합’, 오른쪽 격자의 빈 칸}

\begin{optbox}
초록에 판정어를 쓸 수 있는 사전 등록 대비 10개 가운데 다중 비교 보정 뒤 방향이 확정된 것은 6개, 동등 1개, 보정 전만 유의 1개, 미결정 2개입니다.
\end{optbox}

\slideheading{36}{주장과 근거}{0.8}

주장 다섯 가지입니다. 첫째, 라벨 없는 새 지역에서는 물리식이 출발점입니다. 근거는 물리 유사라벨과 섞은 라벨의 차이 1.6~cm와 기존 지도 CCI v5의 23.8~cm입니다. 둘째, 라벨 몇 개의 이득은 계수 재보정입니다. 재보정 2.5~cm, 추가 ML 이득 확인되지 않음, 사전 등록입니다. 셋째, 수십 개부터는 교차검증 선정이 고정 레시피 이상입니다. 0.9에서 1.3~cm, 캐나다 주도, 사후 분석입니다. 넷째, 수백 개에서 잔차 ML의 이득은 약 0.5~cm이고 그 99\%는 기후 격자 사이 성분입니다. 다섯째, 원천 확대, 배치 규칙, 학습기 교체, 고해상 입력은 이득을 더하지 않았습니다. 그래서 병목은 모델이 아니라 정보입니다.

\pointat{표의 다섯 행, 오른쪽 ‘근거 종류’ 열}

\begin{optbox}
표의 ‘안전’은 비열등 시험 통과가 아니라 큰 실패를 막는 출발점이라는 뜻으로 한정합니다. 다섯째의 근거는 물리식과 제품의 적층, 알래스카에서 고른 배치 규칙의 이전, 학습기 10종, 10에서 500~m 입력 시험이 모두 이득을 더하지 않았다는 결과입니다.
\end{optbox}

\slideheading{37}{한계와 해석}{0.7}

한계 다섯과 대응입니다. 첫째, 학습 원천의 78에서 94\%가 알래스카입니다. 지역 홀드아웃으로 외삽을 평가했지만 확인적 독립 지역이 4에서 5개라는 사실은 남습니다. 둘째, 배치 알고리즘이 무작위로 귀결됐습니다. 분산 배치 이득이 캐나다 −2.5, 레나델타 +1.7~cm로 지역에 따라 다르다는 사실 자체가 지침입니다. 셋째, 격자 안 이득이 0.15~cm로 작습니다. 정보 해상도의 결론입니다. 넷째, 워크플로 검증이 같은 지역 재분할이라 독립 검증이 아니며 새 지역 결과는 10월 11일 뒤 확정됩니다. 다섯째, 지역 안 시험은 사후 분석이므로 사전 등록 시험과 분리해 표시합니다.

\pointat{표의 왼쪽 ‘한계’ 열 다섯 항목}

\slideheading{38}{기대효과와 향후 계획}{0.5}

기대효과는 라벨 투자 판단 근거, 방법 선택 지침, 평가 표준, 지도의 정보 해상도 명시의 넷입니다. 계획은 셋입니다. 격자 안 입력 시험은 WorldCover 10~m 피복 비율과 Sentinel-2 20~m 식생 지수로 먼저 했는데, 격자 안 오차는 0.15~cm 줄어 0.5~cm 기준에 못 미쳤고 레나델타 전체 오차는 1.5~cm 줄었습니다. 위성 적설 자료는 10월 8일까지 들어옵니다. 독립 지역은 극지연구소와 KPDC 자료 협업으로 확보하고, 논문은 Scientific Reports에 제출합니다.

\pointat{오른쪽 표의 ‘격자 안 입력 시험’ 행}

\slideheading{39}{Thank you}{0.2}

이상으로 마치겠습니다. 질문 주시면 답변드리겠습니다. 부록에 학습기별 결과, 판정 단어의 정의, 지도 세부, 추가 실험 요약이 있습니다.
